% Nouveau EGA V (cote 25): archive cote 25, site page 9 (PDF page 10).
% typescript. Draft compared: none.
% AI-transcribed by Claude Opus 5.5 (claude-opus-5-5, adaptive thinking, effort high), Message Batches API, batch ega25b (msgbatch_01A7n5YuXubUPNK1zsepi8RG), 2026-11-13; not yet checked by eye.
\scanpage{25/9}
\mspage{6}

dimension absolue n-1\var{\inserted{$< n = \dim \check{P}$}}{$< n = \dim \check{P}$}{insertion}{ega25_25_009-v0001}, donc le morphisme $\underline{Y}\longrightarrow \overset{\hand{\vee}}{P}$ ne peut être dominant,
donc sa fibre générique $Y_{\hand{\eta}}$ est vide.

\uline{Corollaire} 2.3. Soit $f:X\longrightarrow P$ un morphisme de type fini, et soit Z une
partie constructible de X. Pour que son image inverse dans $Y_{\hand{\eta}}$ soit vide,
il f et s que l'image f(Z) soit finie. En particulier, pour que \var{\xout{X}}{}{typewriter strike-out}{ega25_25_009-v0002} $Y_{\hand{\eta}}$ soit
vide, il f et s que f(X) soit fini.

\uline{Corollaire} 2.4. Soient Z,Z' deux parties fermées \var{\deleted{irréductibles}}{}{deletion}{ega25_25_009-v0003}\var{\inserted{(avec Z irréductible)}}{(avec Z irréductible)}{insertion}{ega25_25_009-v0004} de X,
et soient $Z_{\hand{\eta}}$ et $Z'_{\hand{\eta}}$ leurs images inverses dans $Y_{\hand{\eta}}$ . Pour qu'on ait
$Z_{\hand{\eta}}\subset Z'_{\hand{\eta}}$ , il faut et il suffit que f(Z) soit fini ou que l'on ait $Z\subset Z'$.
Pour que $Z_{\hand{\eta}} = Z'_{\hand{\eta}}$ , il f et s que \var{\xout{ZxxxxZ}}{}{typewriter strike-out}{ega25_25_009-v0005} f(Z) et f(Z') soient finis, ou \var{\inserted{que $Z=Z'$}}{que $Z=Z'$}{insertion}{ega25_25_009-v0006}
\var{\xout{\unsure{En effet, il faut exprimer}}}{}{typewriter strike-out}{ega25_25_009-v0007} C'est une conséquence \var{\xout{\unsure{facile de 2.}}}{}{typewriter strike-out}{ega25_25_009-v0008}
\var{\xout{\unsure{2.3}}}{}{typewriter strike-out}{ega25_25_009-v0009} immédiate de 2.3. , car on voit que $f(Z - Z\cap Z')$ ne peut être fini
que si $Z\subset Z'$, ou si f(Z) est fini (car si on n'a pas $Z\subset Z'$, alors $Z-Z\cap Z'$
est dense dans Z, donc $f(Z-Z\cap Z')$ est dense dans f(Z), et si le premier
est fini donc fermé - étant constructible - il en est de même du second.

\uline{Corollaire} 2.5. A toute composante irréductible $X_i$ de X telle que
$\dim \overline{f(X_i)} > 0$ , faisons correspondre \var{\xout{\illegible{1}}}{}{typewriter strike-out}{ega25_25_009-v0010} son image inverse\var{\xout{s}}{}{typewriter strike-out}{ega25_25_009-v0011} $Y_{i\hand{\eta}}$ dans
$Y_{\hand{\eta}}$ . Alors $Y_{i\hand{\eta}}$ est une composante irréductible de $Y_{\hand{\eta}}$ , et on obtient de
cette façon une correspondance biunivoque entre l'ensemble des comp. irr.
$X_i$ de X telles que $\dim \overline{f(X_i)}> 0$, et l'ens des comp \unsure{\var{\xout{irr}}{}{typewriter strike-out}{ega25_25_009-v0012}|irr} de $Y_{\hand{\eta}}$ .

En effet, il résulte de 2.3. que \var{\xout{X}}{}{typewriter strike-out}{ega25_25_009-v0013} $Y_{\hand{\eta}}$ est réunion des $Y_{i\hand{\eta}}$ \unsure{envisagés|envisagé}
et que ces derniers sont des parties fermées non vides de Y; elles sont
d'ailleurs irréductibles en vertu de 2.1. Enfin, elles sont mutuellement
sans relation d'inclusion en vertu de 2.4. , d'où la conclusion.

Notons que si \var{\xout{\unsure{$d_i$}xxx}}{}{typewriter strike-out}{ega25_25_009-v0014} $\dim X_i$\var{\inserted[typed]{$= d_i$}}{$= d_i$}{insertion}{ega25_25_009-v0015}, \var{\xout{xxxxx}}{}{typewriter strike-out}{ega25_25_009-v0016} on a $\dim Y_i = d_i-1$ , \var{\deleted{comme
il résulte de la proposition}}{}{deletion}{ega25_25_009-v0017} plus générale\hand{ment}:
\uline{Proposition} 2.6. Supposons que \var{\inserted{pour toute comp. irréd. $X_i$ de X on ait}}{pour toute comp. irréd. $X_i$ de X on ait}{insertion}{ega25_25_009-v0018} $\dim \overline{f(X_{\hand{i}})} > 0$ i.e. $Y_{\hand{i\eta}} \neq \emptyset$\var{\inserted{ou que f soit une immers. et $\dim \overline{f(X)}>0$}}{ou que f soit une immers. et $\dim \overline{f(X)}>0$}{insertion}{ega25_25_009-v0019}. Alors on a
\[ \dim Y_{\hand{\eta}} = \dim X - 1 \qquad . \]

\mspage{9}
